\chapter{Conclusiones y Trabajo Futuro}

\section{Síntesis de Cumplimiento de Objetivos}
Las metas delineadas en la fundamentación de este informe, específicamente alineadas con los requisitos del examen final estandarizado, fueron cumplimentadas en su totalidad de manera holística:
\begin{enumerate}
    \item \textbf{Resolución Predictiva:} Se logró modelar la alta varianza del mercado de Nueva York (donde dominaban patrones de ruido y asimetría temporal severa), absorbiendo un 59\% de la varianza total a través de un algoritmo robusto a sobreajustes (\textit{SVR\_RBF}).
    \item \textbf{Ingeniería Híbrida:} La superación del Baseline se fundamentó mayoritariamente por la codificación logarítmica para lidiar con el MAR detectado y la adición del preprocesamiento de coordenadas Haversine, reducidas vectorialmente en un espacio estable vía PCA.
    \item \textbf{Conocimiento Latente:} El *Clustering* dinámico brindó la validación final que las interacciones del mercado de alquiler están determinadas sustancialmente por la exclusividad de uso y la distancia a los focos masivos de consumo (\textit{Times Square}, etc.).
\end{enumerate}

\section{Limitaciones del Estudio}
El rigor científico exige el reconocimiento formal de los límites arquitectónicos presentes:
\begin{itemize}
    \item \textbf{Falta de Variabilidad Estacional:} Al operar exclusivamente sobre una instantánea fotográfica del año 2019, es matemáticamente imposible extraer factores de picos de demanda estival, eventos temporales masivos o la subsecuente pandemia global del COVID-19, dictando que el modelo captura estructuras estructurales puras, pero no variaciones macroeconómicas transitorias.
    \item \textbf{Carencia de Contexto Visual o Descriptivo:} En mercados altamente subjetivos, factores no estructurados como la limpieza general inferida por las fotografías de la galería y la verborragia del *Host* en su descripción textual son predictores ocultos inalcanzables en los datos tabulares disponibles.
\end{itemize}

\section{Recomendaciones y Trabajo Futuro}
Con base en los resultados, el \textit{Action Plan} para futuras iteraciones y para un posible salto hacia la comercialización de la herramienta recae en:
\begin{enumerate}
    \item \textbf{Análisis de Sentimiento Computacional (NLP):} Procesar masivamente el corpus textual de los comentarios empleando modelos de Lenguaje Neuronal masivos (LLMs), agregando un coeficiente de penalización/bonificación basado en polaridad.
    \item \textbf{Motor de Retroalimentación en Tiempo Real:} Diseñar la arquitectura analítica orientada a flujos y *Event Streams*, re-entrenando pesos dinámicos en los modelos SVR con APIs para adaptar las rentabilidades predichas a shocks de la demanda de alta inmediatez logística.
\end{enumerate}
